% Created 2024-01-04 jue 02:01
% Intended LaTeX compiler: pdflatex
\documentclass[11pt]{article}
\usepackage[utf8]{inputenc}
\usepackage[T1]{fontenc}
\usepackage{graphicx}
\usepackage{longtable}
\usepackage{wrapfig}
\usepackage{rotating}
\usepackage[normalem]{ulem}
\usepackage{amsmath}
\usepackage{amssymb}
\usepackage{capt-of}
\usepackage{hyperref}
\usepackage{minted}

\usepackage{parskip}
\usepackage[utf8]{inputenc} %% For unicode chars
\usepackage{placeins}
\usepackage[margin=2.50cm]{geometry}
\usepackage[T1]{fontenc}
\usepackage{mathpazo}
\linespread{1.05}
\usepackage[scaled]{helvet}
\usepackage{courier}
\hypersetup{colorlinks=true,linkcolor=blue}
\RequirePackage{fancyvrb}
\DefineVerbatimEnvironment{verbatim}{Verbatim}{fontsize=\small,formatcom = {\color[rgb]{0.5,0,0}}}
\usepackage{mdframed}
\BeforeBeginEnvironment{minted}{\begin{mdframed}}
\AfterEndEnvironment{minted}{\end{mdframed}}
\setcounter{secnumdepth}{3}
\author{José Manuel Sánchez Álvarez}
\date{}
\title{PHP CRUD II}
\hypersetup{
 pdfauthor={José Manuel Sánchez Álvarez},
 pdftitle={PHP CRUD II},
 pdfkeywords={},
 pdfsubject={},
 pdfcreator={Emacs 28.1 (Org mode 9.6.12)}, 
 pdflang={Spanish}}
\begin{document}

\maketitle
\setcounter{tocdepth}{3}
\tableofcontents

\setlength\parindent{10pt}

\section{CRUD Refactoring}
\label{sec:orga782a31}
Creating a PHP CRUD application often involves repetitive code,
especially for database connections and HTML structures. Using PHP's
\texttt{include}, \texttt{include\_once}, \texttt{require}, and \texttt{require\_once} statements can
streamline this process by allowing you to write common code in separate
files and include them where needed. Here's a guide on how to use these
statements effectively, along with tips for refactoring a CRUD
application.

\subsection{PHP include, include\_once, require, require\_once}
\label{sec:orga1f8f82}
\begin{enumerate}
\item \textbf{include}:
\begin{itemize}
\item Used to include and evaluate a specified file.
\item If the file is not found, PHP issues a warning, but the script will
continue execution.
\end{itemize}

\begin{minted}[]{php}
<?=
include 'filename.php';
\end{minted}

\item \textbf{include\_once}:

\begin{itemize}
\item Similar to \texttt{include}, but it will only include the file once, even
if the statement appears multiple times.
\item Useful to prevent re-declaring functions, classes, or variable
values.
\end{itemize}

\begin{minted}[]{php}
<?=
include_once 'filename.php';
\end{minted}

\item \textbf{require}:

\begin{itemize}
\item Similar to \texttt{include}, but \texttt{require} will produce a fatal error
(halting script execution) if the file cannot be included.
\item Use \texttt{require} when the file is essential for the application to
run.
\end{itemize}

\begin{minted}[]{php}
<?=
require 'filename.php';
\end{minted}

\item \textbf{require\_once}:

\begin{itemize}
\item Combines \texttt{require} and \texttt{include\_once} behavior.
\item Ensures that the file is included only once and execution stops if
the file can't be included.
\end{itemize}

\begin{minted}[]{php}
<?=
require_once 'filename.php';
\end{minted}
\end{enumerate}

\subsection{Refactoring CRUD Application Using Includes}
\label{sec:org87b53b6}
When refactoring your CRUD application, identify parts of your code that
are common across multiple scripts. Typically, these include:

\begin{itemize}
\item Database Connection Setup
\item Header/Footer HTML Code
\item Utility Functions
\end{itemize}

\subsubsection{Example Structure}
\label{sec:orgecb90aa}
\begin{enumerate}
\item \textbf{Database Connection (\texttt{db\_connect.php})}:

\begin{minted}[]{php}
<?=
// db_connect.php
$host = 'localhost';
$db   = 'car_dealership';
$user = 'username';
$pass = 'password';
$charset = 'utf8mb4';

$dsn = "mysql:host=$host;dbname=$db;charset=$charset";
try {
    $pdo = new PDO($dsn, $user, $pass);
    // Set error mode
    $pdo->setAttribute(PDO::ATTR_ERRMODE, PDO::ERRMODE_EXCEPTION);
} catch (\PDOException $e) {
    throw new \PDOException($e->getMessage(), (int)$e->getCode());
}
\end{minted}

\item \textbf{Header/Footer HTML (\texttt{header.php} and \texttt{footer.php})}:

\begin{minted}[]{html}
// header.php
<!DOCTYPE html>
<html>
<head>
    <title>Car Dealership</title>
    <!-- Add other head elements -->
</head>
<body>
\end{minted}

\begin{minted}[]{html}
// footer.php
    <!-- Footer content -->
</body>
</html>
\end{minted}

\item \textbf{Utility Functions (\texttt{utilities.php})}:

\begin{minted}[]{php}
<?=
// utilities.php
function sanitizeInput($data) {
    return htmlspecialchars(stripslashes(trim($data)));
}
\end{minted}
\end{enumerate}

\subsubsection{Implementing Includes in CRUD Scripts}
\label{sec:org6d22398}
In your CRUD scripts, you can now include these files to avoid
redundancy. For example, in a script where you list cars
(\texttt{list\_cars.php}):

\begin{minted}[]{php}
<?php
require_once 'db_connect.php';
include 'header.php';
// ... Rest of your code to list cars ...
include 'footer.php';
\end{minted}

\subsection{Benefits of Using Includes}
\label{sec:orgc253c7f}
\begin{itemize}
\item \textbf{Maintainability}: Changes in common elements like database
configuration or HTML structure can be made in one place.
\item \textbf{Readability}: Reduces clutter in your scripts, making them easier to
read and manage.
\item \textbf{Reusability}: Commonly used code can be reused across different
scripts.
\end{itemize}

\subsection{Best Practices}
\label{sec:org8ec8af9}
\begin{itemize}
\item Use \texttt{require\_once} for critical files like database connections, where
the script should not proceed if the file is missing.
\item Use \texttt{include\_once} for elements like headers and footers where the
inclusion is important but not critical.
\item Regularly review your scripts to identify opportunities for
modularizing your code with includes.
\end{itemize}

By following this guide, you can refactor your PHP CRUD application to
be more efficient, maintainable, and scalable.
\end{document}
